%!TEX program = xelatex
\documentclass[UTF8,10pt,a4paper,fontset=none]{ctexart}

\usepackage[a4paper,top=1.1cm,bottom=1.1cm,left=1.25cm,right=1.25cm]{geometry}
\usepackage{amsmath,amssymb}
\usepackage{booktabs}
\usepackage{array}
\usepackage[most]{tcolorbox}
\usepackage{xcolor}
\usepackage{tikz}
\usetikzlibrary{arrows.meta,positioning}
\usepackage{enumitem}
\usepackage[hidelinks]{hyperref}

\setCJKmainfont{Heiti SC}
\setCJKsansfont{Heiti SC}
\setCJKmonofont{Hiragino Sans GB}

\definecolor{Ink}{HTML}{172033}
\definecolor{Blue}{HTML}{1F4E79}
\definecolor{Green}{HTML}{2A9D8F}
\definecolor{Orange}{HTML}{C86B1F}
\definecolor{SoftBlue}{HTML}{EAF3F8}
\definecolor{SoftGreen}{HTML}{EEF6EA}
\definecolor{SoftOrange}{HTML}{F8EFE4}

\setlength{\parindent}{1.5em}
\setlength{\parskip}{0.08em}
\linespread{1.02}
\setlist[itemize]{leftmargin=1.25em,itemsep=0.04em,topsep=0.05em}
\setlist[enumerate]{leftmargin=1.35em,itemsep=0.04em,topsep=0.05em}
\ctexset{section={format=\large\bfseries,beforeskip=0.45em,afterskip=0.18em}}

\newtcolorbox{note}[2][]{
  enhanced,colback=#2,colframe=Ink!65,boxrule=0.45pt,arc=0.8mm,
  left=1mm,right=1mm,top=0.7mm,bottom=0.7mm,title=#1,fonttitle=\bfseries\small
}
\tikzset{box/.style={draw,rounded corners=1pt,align=center,font=\scriptsize,minimum height=5.8mm,text width=2.05cm},arr/.style={-{Stealth[length=1.8mm]},line width=0.45pt}}

\title{\vspace{-0.9cm}\textbf{第5讲：RNN命名实体识别}\\[-1mm]\normalsize 从单向序列状态到多粒度上下文融合}
\author{高展沛}
\date{\vspace{-0.45cm}2026年6月}

\begin{document}
\maketitle
\vspace{-0.65cm}

\section{课堂任务：RNN-NER 伪代码}
命名实体识别输入句子 $x_1,\ldots,x_n$，输出同长度 BIO 标签 $y_1,\ldots,y_n$。单向 RNN 基线把左侧历史压入状态：
\[
h_t=\tanh(W_xE[x_t]+W_hh_{t-1}+b),\qquad
s_t=W_yh_t+b_y,
\]
\[
p(y_t\mid x_{\le t})=\mathrm{softmax}(s_t),\qquad
\mathcal{L}=\frac{1}{n}\sum_t-\log p(y_t^\star\mid x_{\le t}).
\]

\begin{note}[训练伪代码]{SoftBlue}
\begin{enumerate}
  \item 建 token 词表和 BIO 标签表，如 \texttt{B-PER,I-PER,B-ORG,I-ORG,O}。
  \item 初始化 embedding、RNN 参数和线性分类器。
  \item 对每个句子从左到右递推 $h_t$，由 $s_t$ 预测当前位置标签。
  \item 对所有位置求交叉熵平均，调用反向传播更新参数。
  \item 验证时把预测 BIO 序列转换为实体 span，按实体级 precision/recall/F1 评测。
\end{enumerate}
\end{note}

\begin{note}[解码伪代码]{SoftGreen}
推理阶段逐词取 $\arg\max_y p(y_t)$ 得到初始标签；随后检查 BIO 合法性。若出现 \texttt{O I-ORG} 或 \texttt{B-PER I-LOC} 这类非法转移，可按规则改为 \texttt{B-ORG} 或置为 \texttt{O}。更好的做法是在顶层接 CRF，由 Viterbi 直接搜索全局合法最优序列。
\end{note}

\section{为什么单向 RNN 不够}
NER 不是孤立词分类。实体边界常由左右词共同决定，例如“苹果 发布 新机”中的“苹果”更像组织，而“买了 苹果”中的“苹果”更像普通名词。单向 RNN 在位置 $t$ 只能直接看到 $x_{\le t}$，右侧证据要到未来才进入状态，当前标签判断天然滞后。同时，逐词 softmax 不知道标签之间的结构约束，容易输出局部高分但整体非法的 BIO 序列。

\section{上下文增强方案}
推荐的工程路线是 BiLSTM/GRU + 字符特征 + CRF：
\[
z_t=[e_t^{word};e_t^{char}],\quad
\overrightarrow h_t=RNN_f(z_t,\overrightarrow h_{t-1}),
\]
\[
\overleftarrow h_t=RNN_b(z_t,\overleftarrow h_{t+1}),\quad
o_t=[\overrightarrow h_t;\overleftarrow h_t].
\]
CRF 分数为
\[
score(x,y)=\sum_t s_t[y_t]+\sum_t A[y_{t-1},y_t],
\]
其中 $A$ 学习标签转移规律，Viterbi 负责全局解码。

\begin{center}
\resizebox{0.96\linewidth}{!}{%
\begin{tikzpicture}[node distance=5mm]
\node[box,fill=SoftBlue,draw=Blue] (a) {word/char\\特征};
\node[box,fill=SoftGreen,draw=Green,right=of a] (b) {BiLSTM\\左右上下文};
\node[box,fill=SoftGreen,draw=Green,right=of b] (c) {线性层\\发射分数};
\node[box,fill=SoftOrange,draw=Orange,right=of c] (d) {CRF\\标签结构};
\node[box,fill=SoftOrange,draw=Orange,right=of d] (e) {实体span\\F1评测};
\draw[arr] (a) -- (b);\draw[arr] (b) -- (c);\draw[arr] (c) -- (d);\draw[arr] (d) -- (e);
\end{tikzpicture}}
\end{center}

\section{突破方案：多粒度上下文记忆}
进一步的创新设计是加入文档级实体记忆。模型先抽取高置信实体，保存字符串、类型分布、别名和最近出现位置；后文遇到简称或歧义词时，把记忆向量与句内向量门控融合：
\[
g_t=\sigma(W_g[o_t;m_t]+b_g),\qquad
r_t=g_t\odot o_t+(1-g_t)\odot m_t.
\]
当句内证据强时依赖 $o_t$，当当前句子模糊时调用文档记忆 $m_t$。若数据和算力允许，可把 BiLSTM encoder 替换为中文 BERT/MacBERT，再保留 CRF 作为结构化解码层。

\begin{note}[最终判断]{SoftOrange}
CBOW 学的是“词通常和谁共现”，NER 要判断“词此刻在上下文中扮演什么角色”。因此强 NER 系统必须同时利用字符形态、左右句内语境、标签转移约束和文档级实体历史；这比单向 RNN 更接近真实语言理解。
\end{note}

\end{document}
